\begin{proof}
For fixed $(i,j)$, the events $W=u_iav_j$ over $a$ are disjoint.
Cauchy--Schwarz bounds the expected sum of absolute frequency errors
by $\sqrt{d/n}$. Changing one observation changes that sum by at most
$2/n$. McDiarmid and a union bound over $R^2$ pairs, together with
Hoeffding for the $R^2+2R$ base entries, give a common event of
probability at least $1-\delta$ on which
\begin{equation}
 \|\widehat H-H\|_{\max},\ \|\widehat h-h\|_\infty,
 \|\widehat q-q\|_\infty\le t_n,
 \qquad
 \sum_a\|\widehat M_a-M_a\|_\infty\le R^2b_n.
 \label{II:U:eq:joint-count-event}
\end{equation}
The last inequality sums over both row and column indices. A
per-pair union bound cannot replace this step by exchanging a sum
over letters with a maximum over rows.

Since $\|H^{-1}\|_\infty\le h_0$, the first two sample thresholds give
\begin{align*}
 \|\widehat H^{-1}\|_\infty&\le2h_0,&
 \|\widehat H^{-1}-H^{-1}\|_\infty&\le2Rh_0^2t_n,\\
 \|D_{\widehat q}^{-1}\|_\infty&\le2/q_*,&
 \|D_{\widehat q}^{-1}-D_q^{-1}\|_\infty&\le2t_n/q_*^2.
\end{align*}
Distinct suffixes imply $\|h\|_1,\|\widehat h\|_1\le1$.
For each fixed core prefix, the words $u_iav_j$ over $(a,j)$ are
distinct, so
$\sum_a\|M_a\|_\infty,\sum_a\|\widehat M_a\|_\infty\le R$.
Also $\|D_q\|_\infty,\|D_{\widehat q}\|_\infty\le1$.
Telescoping the three factors in $hH^{-1}D_q$ gives
$\|a^0-\alpha\|_1\le K_\alpha t_n=e_\alpha$.
Telescoping $D_q^{-1}M_aH^{-1}D_q$ and summing over $a$ gives
\[
 \sum_a\|\bar A_a-A_a\|_\infty
 \le v_+v_-\left[
 \frac{2h_0}{q_*}\sum_a\|\widehat M_a-M_a\|_\infty+K_0t_n\right]
 \le e_A.
\]
The remaining thresholds imply $e_\alpha\le\gamma/4$ and
$e_A\le\min(\gamma,s)/2$. Thus $a^0$ is positive, and its sum is
within $e_\alpha$ of one. Its normalized entries are at least
$(\gamma-e_\alpha)/(1+e_\alpha)\ge\gamma/2$;
$\|\bar\alpha-\alpha\|_1\le2e_\alpha/(1-e_\alpha)$.
The transition entries are at least $\gamma/2$ and the recomputed
stopping probabilities at least $s/2$. All exact checks therefore pass.

Downward rounding changes at most $dR$ entries in each transition row,
with total change at most $dRh_b$; STOP absorbs this mass. Initial TV
changes by at most $(R-1)h_b$. The original source has occupation
sum at most $1/s$. Apply Lemma~\ref{II:U:lem:stopped-coupling}, or
\eqref{II:U:eq:pathwise-general} with the row-remainder identity, to obtain
\eqref{II:U:eq:fixed-basis-tv}. The rounding restriction leaves initial
and transition entries at least $\gamma/4$ and never reduces STOP.

For complexity, a trie of the $O(dR^2)$ queried words and prefixes
permits counting in $O(nD+dR^2D)$ symbol operations. Counts have
$O(\log n)$ bits. Fraction-free elimination computes the rational
inverses; determinants of the integer matrices have polynomial bit
length in $R,\log n$ and the input bit length of $V$. The remaining
$O(dR^3)$ dense rational arithmetic operations, normalization, and
integer floor divisions also have polynomial bit cost. This proves
the stated polynomial bound; it does not assert a small uniform
degree or a fixed-word-size implementation. Dyadic probabilities
are stored as their integer numerators, with remainders determined
exactly, giving the output bound.
\end{proof}
